\chapter{Arquitetar uma aplicação Quarkus}
\label{cap:arquitetura-quarkus}

\proximasaulas

A seção sobre \emph{Clean Architecture} do capítulo~\ref{cap:profissional}
terminou com uma leitura da Dependency Rule para Quarkus: anotações de
framework nos adaptadores, regras de negócio em classes comuns. Este
capítulo transforma isso num projeto concreto, com nomes de pacotes,
e mostra onde ficam REST, Panache, configuração, testes, resiliência e
observabilidade.

O exemplo é um serviço de pedidos, com raiz no pacote
\texttt{br.com.turma.pedidos}. As listagens são esqueletos ilustrativos,
escritos para este capítulo; não estão em \texttt{examples/} e não
foram compilados contra uma versão específica do Quarkus. A divisão em
pacotes é uma aplicação do livro de Martin e do artigo de Cockburn, não
uma regra do Quarkus, que não impõe estrutura nenhuma.

\section{Os círculos da Clean Architecture em pacotes Quarkus}
\label{sec:circulos-quarkus}

O capítulo~22 do \citetitle{livro:martin2017} desenha quatro círculos
concêntricos e uma regra: as dependências de código-fonte só apontam
para dentro. O círculo de dentro não sabe o nome de nada que está fora
dele, nem de classes, nem de funções, nem de frameworks. A tabela mapeia
cada círculo para pacotes e construções do Quarkus.

\small
\begin{longtable}{@{}p{2.4cm} p{4.4cm} p{5.4cm}@{}}
\toprule
Círculo & Pacote & O que tem e o que pode importar \\
\midrule \endhead
\emph{Entities} & \texttt{\ldots pedidos.dominio} & Records e classes
  com as regras de negócio (\texttt{Pedido}, \texttt{ItemPedido},
  \texttt{Dinheiro}). Só importa \texttt{java.*}. Nada de Jakarta, nada
  de Quarkus \\
\emph{Use Cases} & \texttt{\ldots pedidos.aplicacao} e
  \texttt{\ldots aplicacao.porta} & Casos de uso (\texttt{FecharPedido})
  e as \emph{portas}: interfaces que o caso de uso precisa
  (\texttt{Pedidos}, \texttt{CobrancaGateway}). Importa só
  \texttt{dominio} \\
\emph{Interface Adapters} & \texttt{\ldots adaptador.\allowbreak entrada.\allowbreak rest},
  \texttt{\ldots adaptador.\allowbreak saida.\allowbreak persistencia},
  \texttt{\ldots adaptador.\allowbreak saida.\allowbreak cobranca} & Resources REST (\texttt{@Path}),
  repositórios Panache e entidades JPA, clientes REST. Implementam as
  portas e convertem entre o modelo externo e o domínio \\
\emph{Frameworks \& Drivers} & \texttt{\ldots config} e o próprio
  Quarkus & \texttt{application.properties}, \texttt{@ConfigMapping},
  producers CDI que montam os casos de uso, extensões, banco, Vert.x \\
\bottomrule
\end{longtable}
\normalsize

\paragraph{Onde ficam as extensões.}
Extensões Quarkus são dependências de build, não pacotes seus: entram no
\texttt{pom.xml} e aparecem no código só nos adaptadores e em
\texttt{config}. Uma forma de conferir é olhar os \texttt{import} de
\texttt{dominio} e \texttt{aplicacao}: se aparecer \texttt{io.quarkus},
\texttt{jakarta.persistence} ou \texttt{jakarta.ws.rs}, a regra foi
quebrada. Os adaptadores, por sua vez, importam exatamente esses
pacotes: um resource com \texttt{@Path} ou uma entidade com
\texttt{@Entity} junta, na mesma classe, os dois círculos de fora. O
\citetitle{livro:martin2017} descreve o círculo externo como, em geral,
pouco código além da \emph{cola} que fala com o círculo seguinte, e as
anotações são essa cola. Aceitamos a fusão porque a seta continua
apontando para dentro: o que a regra precisa proteger são os dois
círculos internos, e eles não importam nada do framework.

\paragraph{E a injeção de dependência?}
A pergunta honesta é se \texttt{@ApplicationScoped} e \texttt{@Inject}
podem aparecer em \texttt{aplicacao}. São anotações da especificação
Jakarta (\texttt{jakarta.enterprise.context}, \texttt{jakarta.inject}),
não do Quarkus, e não mudam o comportamento da classe num teste JUnit
puro. Há duas escolhas defensáveis, e a seção seguinte mostra as duas:
\begin{description}
\item[Caso de uso anotado] Mais simples; aceita uma dependência fraca de
  uma especificação estável.
\item[Caso de uso sem anotação, criado por um producer em \texttt{config}]
  Mantém \texttt{aplicacao} com zero dependências externas, ao custo de
  um arquivo a mais. Lembre da seção~\ref{sec:jandex}: classe sem
  \emph{bean defining annotation} não é descoberta, então o producer é
  obrigatório.
\end{description}

O que não é defensável é deixar o domínio depender do framework para
existir: uma entidade de domínio que só funciona dentro de uma transação
JPA, ou uma regra de negócio que lê \texttt{ConfigProvider} diretamente.

\paragraph{Configuração como detalhe.}
O guia \citetitle{doc:quarkus-config-mappings} mostra o
\texttt{@ConfigMapping}: uma interface pública cujos métodos viram
propriedades a partir de um prefixo. Ela fica em \texttt{config}; o caso
de uso recebe só os valores de que precisa (um \texttt{Duration}, um
limite), pelo construtor, sem saber de onde vieram.

\section{Portas e adaptadores ligados por CDI}
\label{sec:hexagonal-cdi}

A arquitetura hexagonal de Alistair Cockburn
(\citetitle{artigo:cockburn-hexagonal}, 2005) chega à mesma regra por
outro caminho. A intenção declarada é permitir que a aplicação seja
\enquote{equally be driven by users, programs, automated test or batch
scripts} e desenvolvida e testada isolada dos dispositivos e bancos que
vai usar. Uma \termo{porta}{\emph{Port}, no vocabulário de
\emph{ports \& adapters}: o ponto onde a aplicação conversa com o mundo
de fora, definido como uma API.} identifica uma conversa com propósito;
um \termo{adaptador}{\emph{Adapter}: converte a API da porta nos sinais
de um dispositivo concreto (HTTP, SQL, fila) e vice-versa.} traduz essa
conversa para uma tecnologia. No Quarkus, o CDI é quem liga uma ponta à
outra.

\paragraph{A porta e o caso de uso.}
Em \texttt{aplicacao.porta}, uma interface; em \texttt{aplicacao}, o
caso de uso que depende dela pelo construtor, sem anotação:

\begin{lstlisting}[language=Java]
package br.com.turma.pedidos.aplicacao.porta;
public interface Pedidos {                 // porta de saída
    Optional<Pedido> porId(PedidoId id);
    void salvar(Pedido pedido);
}

package br.com.turma.pedidos.aplicacao;
public final class FecharPedido {          // caso de uso, Java puro
    private final Pedidos pedidos;
    private final CobrancaGateway cobranca;
    public FecharPedido(Pedidos pedidos, CobrancaGateway cobranca) {
        this.pedidos = pedidos; this.cobranca = cobranca;
    }
    public Pedido executar(PedidoId id) { /* regras em dominio */ }
}
\end{lstlisting}

\paragraph{O adaptador e o producer.}
Em \texttt{adaptador.saida.persistencia}, o adaptador é um bean comum que
implementa a porta. Em \texttt{config}, um producer cria o caso de uso:

\begin{lstlisting}[language=Java]
package br.com.turma.pedidos.adaptador.saida.persistencia;
@ApplicationScoped
class PedidosPanache implements Pedidos { /* usa PedidoRepository */ }

package br.com.turma.pedidos.config;
@ApplicationScoped
class CasosDeUso {
    @Produces @Singleton
    FecharPedido fecharPedido(Pedidos p, CobrancaGateway c) {
        return new FecharPedido(p, c);
    }
}
\end{lstlisting}

\paragraph{Por que \texttt{@Singleton} no producer.}
O guia \citetitle{doc:quarkus-cdi} recomenda \texttt{@ApplicationScoped}
por padrão, a menos que haja um bom motivo para \texttt{@Singleton}. A
diferença é o \termo{client proxy}{Objeto gerado que fica no lugar do bean
injetado e delega para a instância real, criada sob demanda.}: um bean
\texttt{@Singleton} não tem proxy. Aqui há um motivo: \texttt{FecharPedido}
é \texttt{final}, e o guia \citetitle{doc:quarkus-cdi-reference} explica
que classes \texttt{final} não podem ser \enquote{proxiadas}, a não ser que
o Quarkus transforme o bytecode removendo o \texttt{final}. Um caso de uso
imutável e sem estado não ganha nada com o proxy. O preço, segundo o
mesmo guia: um bean \texttt{@Singleton} não pode ser substituído por
\texttt{QuarkusMock}; num \texttt{@QuarkusTest} use \texttt{@InjectMock}
com \texttt{@MockitoConfig(convertScopes = true)}.

\paragraph{Por que isso combina com o build time.}
Toda a ligação é estática e visível no índice: o tipo da porta, o bean
que a implementa e o producer. O guia de CDI diz que exatamente um bean
deve ser atribuível a cada ponto de injeção, senão \enquote{the build
fails} (com \texttt{UnsatisfiedResolutionException} ou
\texttt{AmbiguousResolutionException}). Faltou um adaptador ou sobrou um:
você descobre no build, não na primeira requisição. Note também que nada é
\texttt{private} onde o container precisa entrar: injeção pelo construtor
(ou pelo parâmetro do producer), como recomenda o guia
\citetitle{doc:quarkus-cdi-reference} e como a
seção~\ref{sec:gizmo} explicou em bytecode.

\paragraph{Trade-offs.}
\begin{itemize}
\item O producer é código a mais para manter. Num serviço pequeno,
  anotar \texttt{FecharPedido} com \texttt{@ApplicationScoped} é uma
  troca razoável; registre a decisão.
\item Uma porta por adaptador pode virar burocracia. Crie portas onde
  existe uma segunda implementação real ou provável: o dublê de teste
  (Parte~\ref{parte:testes}) já conta como segunda implementação.
\item Separar domínio e entidade JPA significa mapear um para o outro.
  A próxima seção discute quando isso vale.
\end{itemize}

\section{REST e Panache como adaptadores}
\label{sec:panache-rest}

\subsection*{O resource REST é um adaptador de entrada}

Em \texttt{adaptador.entrada.rest}, o resource converte HTTP em chamada
de caso de uso e o resultado em DTO. Ele não tem regra de negócio:

\begin{lstlisting}[language=Java]
package br.com.turma.pedidos.adaptador.entrada.rest;
@Path("/pedidos")
public class PedidoResource {
    private final FecharPedido fecharPedido;
    @Inject
    PedidoResource(FecharPedido fecharPedido) { this.fecharPedido = fecharPedido; }

    @POST @Path("/{id}/fechamento")
    public PedidoDto fechar(@PathParam("id") String id) {
        return PedidoDto.de(fecharPedido.executar(new PedidoId(id)));
    }
}
\end{lstlisting}

\paragraph{Em que thread isso roda.}
O guia \citetitle{doc:quarkus-rest} diz que a thread depende da
assinatura do método. Métodos que devolvem \texttt{Uni},
\texttt{Multi}, \texttt{CompletionStage} ou \texttt{Publisher} são
considerados não bloqueantes e rodam na thread de I/O; os demais rodam
numa \emph{worker thread}. \texttt{@Blocking} e \texttt{@NonBlocking}
mudam o padrão, e um método com \texttt{@Transactional} é tratado como
bloqueante. O resource acima devolve um DTO comum, então roda numa worker
thread, e o caso de uso pode bloquear em JDBC sem travar o servidor. O
capítulo~\ref{cap:decisoes-quarkus} compara esse modelo com virtual
threads e com o reativo.

\subsection*{Panache: active record ou repository}

O guia \citetitle{doc:quarkus-panache} oferece dois estilos:
\begin{description}
\item[Active record] A entidade estende \texttt{PanacheEntity} e as
  consultas são métodos estáticos nela (\texttt{Pedido.findById(id)});
  o guia chama isso de \enquote{the Panache Active Record way}.
\item[Repository] Uma classe \texttt{@ApplicationScoped} implementa
  \texttt{PanacheRepository<PedidoEntity>} e oferece os mesmos métodos;
  estender \texttt{PanacheEntity} fica opcional.
\end{description}

O guia apresenta os dois sem recomendar um. Do ponto de vista da
Dependency Rule, a diferença é grande:

\begin{itemize}
\item \textbf{Active record coloca persistência dentro da entidade.} Se
  essa entidade for o seu \texttt{Pedido} de \texttt{dominio}, o círculo
  de dentro passa a importar \texttt{io.quarkus.hibernate.orm.panache} e
  \texttt{jakarta.persistence}, e uma regra de negócio pode chamar
  \texttt{persist()} no meio de um cálculo. A regra quebra no lugar mais
  central.
\item \textbf{Repository cabe num adaptador.} \texttt{PedidoEntity} (JPA)
  e \texttt{PedidoRepository} (Panache) ficam em
  \texttt{adaptador.saida.persistencia}; \texttt{PedidosPanache}
  implementa a porta \texttt{Pedidos} e converte entre
  \texttt{PedidoEntity} e o record \texttt{Pedido}.
\end{itemize}

\paragraph{O custo honesto.}
O mapeamento entidade$\leftrightarrow$domínio é código repetitivo, e
para um serviço que é essencialmente CRUD ele não protege nada: não há
regra de negócio para isolar. Nesse caso, active record é a escolha mais
simples e não é \enquote{errado}; é uma decisão consciente de que o
banco \emph{é} o modelo. Quando aparecer a primeira regra que não cabe
num \texttt{UPDATE}, extraia o domínio. Registre a escolha num ADR para
que o próximo dev saiba que foi intencional.

\paragraph{Testabilidade, de novo.}
O guia de Panache aponta mais uma diferença: com active record, o
Mockito não consegue simular os métodos estáticos, e é preciso usar o
\texttt{PanacheMock}; com repository, \enquote{you can use Mockito
directly}. Com a porta \texttt{Pedidos} no meio, nem isso: o teste do
caso de uso usa um fake em memória, como na Parte~\ref{parte:testes}.

\paragraph{O que o Panache faz no seu bytecode.}
Os getters e setters das entidades com campos \texttt{public} são
gerados no build (seção~\ref{sec:gizmo}). O guia descreve esse
\emph{enhancement} para as entidades; os records do domínio não são
entidades JPA, então continuam exatamente como o \texttt{javac} os
gerou (seção~\ref{sec:records}).

\section{Testes por camada}
\label{sec:testes-camadas}

A arquitetura acima existe, em boa parte, para tornar os testes da
Parte~\ref{parte:testes} baratos. O capítulo~\ref{cap:ponte-quarkus} apresentou
\texttt{@QuarkusTest}, \texttt{@QuarkusIntegrationTest} e Dev Services;
aqui a pergunta é \emph{onde} usar cada um.

\small
\begin{longtable}{@{}p{3cm} p{4.2cm} p{5cm}@{}}
\toprule
Camada & Tipo de teste & Por quê \\
\midrule \endhead
\texttt{dominio} & JUnit puro, rápido, muitos & Sem framework, sem
  container; o ciclo de TDD da aula~1 roda em milissegundos \\
\texttt{aplicacao} & JUnit puro com fakes das portas & As portas são
  interfaces suas; um \texttt{Pedidos} em memória substitui o banco \\
\texttt{adaptador.\allowbreak saida.*} & \texttt{@QuarkusTest} com Dev Services &
  O que importa é o SQL, o mapeamento e a configuração reais \\
\texttt{adaptador.\allowbreak entrada.\allowbreak rest} & \texttt{@QuarkusTest} pela API HTTP,
  com o caso de uso substituído quando convém & Serialização, códigos de
  status e validação \\
Artefato final & \texttt{@QuarkusIntegrationTest} & Caixa preta contra o
  JAR ou o executável nativo produzido pelo build \\
\bottomrule
\end{longtable}
\normalsize

\paragraph{Substituir um bean num \texttt{@QuarkusTest}.}
O guia \citetitle{doc:quarkus-testing} mostra o \texttt{@InjectMock},
disponível com a dependência \texttt{quarkus-junit-mockito}, e o
\texttt{QuarkusMock}, com \texttt{installMockForType} para instalar um
mock antes dos testes. Na arquitetura hexagonal, o candidato natural a
mock é uma porta (\texttt{CobrancaGateway}), nunca uma classe de
domínio. O \texttt{QuarkusMock} só funciona com beans de escopo normal;
para o \texttt{FecharPedido} \texttt{@Singleton}, use \texttt{@InjectMock}
com \texttt{@MockitoConfig(convertScopes = true)}.

\paragraph{Cuidado com o custo de subir o Quarkus.}
O mesmo guia diz que, nos exemplos dele, o Quarkus sobe uma vez para
todos os testes. Um teste com \texttt{@TestProfile} diferente do anterior
faz o Quarkus ser desligado e iniciado de novo com o novo perfil. Agrupe
os testes por perfil e prefira poucos perfis: cada um soma uma subida ao
tempo da suíte.

\paragraph{Dev Services e o modo nativo.}
Os Dev Services (\citetitle{doc:quarkus-dev-services}) sobem banco e
outros serviços com Testcontainers quando nada foi configurado, nos
modos de desenvolvimento e teste. O guia de testes diz que o
\texttt{@QuarkusIntegrationTest} funciona com contêineres iniciados por
Dev Services, o que permite rodar os mesmos testes de caixa preta contra
o executável nativo no CI. É o teste que pega os problemas da
seção~\ref{sec:reflexao-nativo}: reflexão não registrada só aparece no
binário nativo.

\paragraph{Mutação.}
O PIT (aula~2) brilha em \texttt{dominio} e \texttt{aplicacao}, onde os
testes são rápidos e o código não depende do Quarkus. A integração do
PIT com \texttt{@QuarkusTest} continua não verificada
(capítulo~\ref{cap:ponte-quarkus}); manter as regras fora do framework
evita depender dela.

\section{Resiliência nos adaptadores de saída}
\label{sec:resiliencia-quarkus}

A aula~2 e o capítulo~\ref{cap:ponte-quarkus} apresentaram retry,
timeout, fallback, circuit breaker e bulkhead, e o SmallRye Fault
Tolerance que o Quarkus usa. A pergunta arquitetural é \emph{onde}
colocar as anotações \texttt{@Retry}, \texttt{@Timeout},
\texttt{@Fallback}, \texttt{@CircuitBreaker} e \texttt{@Bulkhead} que o
guia \citetitle{doc:quarkus-fault-tolerance} mostra.

\paragraph{Na borda, perto da chamada que falha.}
A falha de rede é um detalhe do adaptador: o domínio não sabe que a
cobrança é HTTP. As anotações vão no adaptador de saída que implementa a
porta:

\begin{lstlisting}[language=Java]
package br.com.turma.pedidos.adaptador.saida.cobranca;
@ApplicationScoped
class CobrancaHttp implements CobrancaGateway {
    @Inject @RestClient CobrancaApi api;   // package-private, sem reflexão

    @Timeout(800) @Retry(maxRetries = 2) @CircuitBreaker
    public ResultadoCobranca cobrar(Pedido p) { /* converte e chama api */ }
}
\end{lstlisting}

\paragraph{O que a porta devolve quando tudo falha.}
O \texttt{@Fallback} é uma decisão de negócio disfarçada de técnica:
\enquote{cobrança indisponível} vira pedido pendente, erro para o usuário
ou fila? Exponha esse resultado no tipo de retorno da porta
(\texttt{ResultadoCobranca} com um caso \texttt{Indisponivel}, por
exemplo, com um \texttt{sealed interface} da
seção~\ref{sec:sealed-pattern}), para que o caso de uso decida e o teste
unitário cubra a decisão.

\paragraph{Bulkhead e virtual threads.}
O bulkhead limita execuções concorrentes. Com virtual threads, a regra
da seção~\ref{sec:virtual-threads} é a mesma: não limite com pool de
threads, limite o recurso escasso. O bulkhead no adaptador é o lugar
natural para isso.

\paragraph{Testando.}
Os padrões da aula~2 se testam no adaptador: um \texttt{@QuarkusTest}
com o serviço externo simulado (por exemplo, com o Toxiproxy da aula~2)
mostra se timeout e circuit breaker se comportam como configurado. O
caso de uso, por sua vez, testa só a reação a
\texttt{Indisponivel}, com um fake da porta.

\section{Observabilidade sem sujar o domínio}
\label{sec:observabilidade-quarkus}

Observabilidade é outro detalhe técnico: a regra de fechar pedido não
muda porque alguém quer um gráfico. As três peças abaixo entram por
extensão e ficam nos adaptadores e em \texttt{config}.

\paragraph{Health checks.}
O guia \citetitle{doc:quarkus-health} expõe \texttt{/q/health/live}
(\enquote{The application is up and running}),
\texttt{/q/health/ready} (pronta para atender), \texttt{/q/health/started}
e \texttt{/q/health}, que agrega todos. Um check é um bean que implementa
\texttt{HealthCheck} anotado com \texttt{@Liveness} ou \texttt{@Readiness}.
Extensões registram checks sozinhas: o guia cita a
\texttt{quarkus-agroal}, que registra um readiness check para cada
datasource. Coloque checks próprios em \texttt{adaptador} (por exemplo,
\texttt{adaptador.saida.cobranca.CobrancaHealth}), perto do recurso que
eles verificam.

Liveness e readiness respondem a perguntas diferentes, e confundi-las
custa caro num orquestrador: um liveness que depende do banco faz o
container ser reiniciado quando o banco cai, o que não conserta o
banco. Use dependências externas no readiness.

\paragraph{Métricas.}
O guia \citetitle{doc:quarkus-micrometer} diz que \enquote{Micrometer is
the recommended approach to metrics for Quarkus}; com a extensão
\texttt{micrometer-registry-prometheus}, as métricas saem em
\texttt{/q/metrics}. As anotações \texttt{@Timed} e \texttt{@Counted}
funcionam em métodos de beans. Ponha-as no resource ou no adaptador;
se uma métrica for de negócio (pedidos fechados por hora), publique um
evento de domínio e conte-o num adaptador, em vez de importar Micrometer
em \texttt{dominio}.

\paragraph{Traces.}
O guia \citetitle{doc:quarkus-opentelemetry} diz que, com a dependência
adicionada, a extensão gera dados de tracing por padrão e os exporta em
lote por OTLP/gRPC para \texttt{http://localhost:4317}. Métricas e logs
via OpenTelemetry são \emph{tech preview} e vêm desligados; ligam-se em
build time com \texttt{quarkus.otel.metrics.enabled=true} e
\texttt{quarkus.otel.logs.enabled=true}.

\paragraph{E a JVM?}
Essas três peças mostram a aplicação. Para ver a JVM (pausas de GC,
alocação, pinning de virtual threads, compilação), as ferramentas são as
da seção~\ref{sec:ferramentas}: JFR, \texttt{jcmd} e JMC. Os dois níveis
se completam: o trace mostra \emph{qual} requisição ficou lenta; o JFR
mostra \emph{por quê}, por exemplo uma pausa de GC no mesmo instante.

\section{Onde aprofundar}

\subsection*{Livros e artigos}
\small
\begin{longtable}{@{}p{4.4cm} p{3cm} p{4.8cm}@{}}
\toprule
Fonte & Autor(es), ano & Por que ler \\
\midrule \endhead
\citetitle{livro:martin2017} \pt{} (\emph{Arquitetura Limpa})
  & \citeauthor{livro:martin2017}, \citeyear{livro:martin2017}
  & Caps.~22 (os círculos e a Dependency Rule) e 32 (o framework como
  detalhe) \\
\citetitle{artigo:cockburn-hexagonal}
  & \citeauthor{artigo:cockburn-hexagonal}, 2005
  & O texto original de ports \& adapters, com a motivação de testar a
  aplicação sem UI nem banco \\
\bottomrule
\end{longtable}
\normalsize

\subsection*{Documentação}
\citetitle{doc:quarkus-cdi}, \citetitle{doc:quarkus-cdi-reference},
\citetitle{doc:quarkus-rest}, \citetitle{doc:quarkus-panache},
\citetitle{doc:quarkus-config-mappings},
\citetitle{doc:quarkus-testing}, \citetitle{doc:quarkus-dev-services},
\citetitle{doc:quarkus-fault-tolerance}, \citetitle{doc:quarkus-health},
\citetitle{doc:quarkus-micrometer} e
\citetitle{doc:quarkus-opentelemetry}.

\subsection*{No IDE: blog da JetBrains}
Outros posts da JetBrains no apêndice~\ref{ap:jetbrains}.
\begin{itemize}
\obra*{artigo:jetbrains-testcontainers-spring-boot}{Testes de integração
  com Testcontainers; o exemplo é Spring Boot, mas a ideia é a mesma dos Dev
  Services do Quarkus.}
\obra*{curso:jetbrains-guide-javax-jakarta}{Migrar de \texttt{javax} para
  \texttt{jakarta}, o namespace que o Quarkus usa.}
\end{itemize}

\subsection*{Pesquisa acadêmica}
Não encontramos estudo acadêmico revisado por pares sobre arquitetura
hexagonal ou Clean Architecture; os itens abaixo são a base conceitual e
empírica por trás delas. Acesso e links no apêndice~\ref{ap:pesquisa}.
\begin{itemize}
\pesquisa{paper:parnas-modules-1972}{A raiz de portas e adaptadores:
  esconder atrás de uma interface o que pode mudar.}
\pesquisa{paper:garlan-shaw-intro-softarch-1994}{Relatório técnico da CMU
  sobre estilos arquiteturais: camadas, eventos, pipes and filters.}
\pesquisa{paper:perry-wolf-foundations-1992}{Define arquitetura e introduz
  as ideias de erosão e desvio arquitetural.}
\pesquisa{paper:desilva-erosion-survey-2012}{Técnicas contra a erosão, como
  regras de dependência no estilo do ArchUnit.}
\pesquisa{paper:dragoni-microservices-2017}{Panorama histórico de
  microsserviços e dos problemas em aberto.}
\pesquisa{paper:soldani-microservices-pains-gains-2018}{O que praticantes
  relatam de bom e de ruim em microsserviços.}
\pesquisa{paper:viggiato-microservices-practice-2018}{Pesquisa da UFMG com
  122 profissionais sobre microsserviços na prática.}
\pesquisa{paper:spadini-mock-objects-java-2019}{Por que se mocka em Java e
  como os mocks evoluem com o código.}
\pesquisa{paper:almeida-mock-objects-sbqs-2023}{Da UFF: Mockito é o
  framework mais usado em Java e dependências externas são as mais mockadas.}
\pesquisa*{livro:valente2020}{Cap.~7 (arquitetura: camadas, MVC,
  microsserviços, mensageria) e cap.~8 (testes).}
\end{itemize}

\section{Desafios práticos}

\desafio{Pegue um serviço Quarkus que você já tem e liste os
\texttt{import} de cada pacote. Desenhe as setas de dependência entre
pacotes. Alguma aponta para fora, do domínio para o framework? Proponha
a menor mudança que inverte essa seta.}

\desafio{Implemente o caso de uso \texttt{FecharPedido} da
seção~\ref{sec:hexagonal-cdi} em TDD, com um fake em memória da porta
\texttt{Pedidos}, sem nenhuma dependência do Quarkus no módulo. Quanto
tempo leva a suíte? Compare com a mesma regra testada por
\texttt{@QuarkusTest}.}

\desafio{Escreva a mesma entidade nos dois estilos do Panache (active
record e repository). Para cada um, responda: onde fica a regra de
negócio, como se testa sem banco e o que muda se amanhã a persistência
for um serviço HTTP em vez de um banco.}

\desafio{No adaptador de cobrança, configure \texttt{@Timeout} e
\texttt{@CircuitBreaker} e use o Toxiproxy (aula~2) para injetar
latência. Verifique no trace do OpenTelemetry e nas métricas do
Micrometer o que o usuário viu enquanto o circuito estava aberto.}

\desafio{Escreva um readiness check para a cobrança e decida, com
argumentos, se ele deve também ser liveness. O que o orquestrador faria
em cada caso quando a cobrança cai?}

